\conclusion

В ходе работы разработана и исследована параллельная версия программного комплекса HeatSim3D, предназначенного для численного моделирования нестационарной теплопередачи в трёхмерных сплошных средах. Распараллеливание выполнено средствами технологии OpenMP для многоядерных процессоров с общей памятью. Работа опирается на программный комплекс, созданный в предыдущем семестре, и сохраняет его математическую модель, численный метод и графический интерфейс.

В первой главе кратко изложена математическая модель на основе уравнения Фурье--Кирхгофа с граничными условиями первого рода на шести гранях и явная конечно-разностная схема FTCS с условием устойчивости. Рассмотрены модель исполнения OpenMP, способы распределения итераций между потоками и механизм редукции. Показано, что ускорение параллельной программы ограничивается последовательными участками в соответствии с законом Амдала, а для стенсильных вычислений --- также пропускной способностью памяти.

Во второй главе на основе анализа информационных зависимостей установлено, что все внутренние узлы на одном шаге по времени можно обрабатывать независимо, если хранить текущий и новый слои температурного поля в разных массивах. Распараллелены основной цикл схемы, наложение граничных условий и вычисление средней температуры. Ручное слияние двух внешних циклов обеспечило равномерную загрузку потоков и совместимость со стандартом OpenMP 2.0, который поддерживается компилятором Microsoft Visual C++. В расчётный модуль добавлено управление числом потоков, а программный комплекс собирается и работает в операционных системах Linux и Windows.

В третьей главе подтверждена корректность параллельной версии. Максимальное отклонение численного решения от аналитического на сетке $31^3$ составило 0,20~°C, или 0,25~\% от перепада температур, а при измельчении сетки погрешность уменьшается в четыре раза при каждом уменьшении шага вдвое, что соответствует второму порядку сходимости. Результаты расчёта при одном и четырёх потоках, а также в версиях для Linux и Windows совпадают побитно. Исследование производительности на четырёхъядерном процессоре показало, что для сеток от миллиона узлов ускорение на четырёх потоках составляет 3,2--3,4 раза при эффективности 0,81--0,86. На малой сетке $50^3$ эффективность снижается до 0,68 из-за относительно большой доли затрат на синхронизацию потоков. Частая последовательная запись результатов способна снизить ускорение почти втрое, что согласуется с законом Амдала и позволяет дать практическую рекомендацию по выбору интервала сохранения.

Таким образом, цель работы достигнута, и параллельная версия позволяет выполнять расчёты на мелких сетках в несколько раз быстрее при неизменной точности. Направлениями дальнейшего развития являются асинхронная запись результатов в отдельном потоке, блочная обработка расчётной области для лучшего использования кэша процессора и исследование масштабируемости на процессорах с большим числом ядер.

\fillme{Перечень компетенций ниже повторяет список прошлого семестра. Обязательно сверить с учебным планом седьмого семестра: при несоответствии работа снимается с защиты.}

Перечень компетенций для направления 09.03.04 <<Программная инженерия>>, сформированных в ходе производственной практики (научно-исследовательской работы):

УК-1. Способен осуществлять поиск, критический анализ и синтез информации, применять системный подход для решения задач. Компетенция сформирована при изучении литературы по технологии OpenMP и показателям эффективности параллельных программ, а также при анализе ограничений, связанных с законом Амдала и пропускной способностью памяти.

УК-4. Способен осуществлять деловую коммуникацию в устной и письменной формах на государственном языке Российской Федерации и иностранном(ых) языке(ах). Компетенция проявилась при работе с англоязычными источниками по параллельному программированию, подготовке отчётной документации и обсуждении результатов с руководителем практики.

ПК-1. Способен проводить научно-исследовательские и опытно-конструкторские разработки. Компетенция сформирована при разработке параллельной версии расчётного модуля, планировании и проведении вычислительных экспериментов по оценке ускорения и эффективности и анализе их результатов.

ПК-4. Способен создавать и анализировать требования на разработку программно-информационных систем. Компетенция проявилась при формулировании требований к параллельной версии, включая совпадение результатов с последовательной версией, управление числом потоков и работу в операционных системах Linux и Windows.

ПК-5. Способен осуществлять концептуальное, функциональное и логическое проектирование программно-информационных систем. Компетенция проявилась при проектировании организации данных для параллельной обработки, выборе схемы распределения вычислений между потоками и разработке средств автоматизации вычислительных экспериментов.
